\section*{About this document}
\markboth{About this document}{About this document}

This document follows \emph{Groundwater Audit from Space} (the v2 report and its experiments log). That work
asked whether satellites can find and size single pumping wells; the answer for Semarang was no. This
document designs the app that uses what did work. Its research basis is experiment file 15 and files
04--14 of the experiments log \cite{exp15}.

\textbf{Assumptions.} Study area: Central Java for module A, Semarang--Demak for B and C. Team: four
students. Time: six weeks to the national round. Data: only open data and documents a citizen can request.

\subsection*{How to read this document}

Every item that matters has a tag: a letter for its kind and a circled number. For example, \Q{1} (how much
is pumped and paid?) is answered by method \M{1}, checked by \V{1}, built in component \Sx{2}, and
delivered as output \Ox{1}. Section~\ref{sec:trace} lists every link.

\begin{center}
\small\renewcommand{\arraystretch}{1.0}
\begin{tabular}{@{}llp{62mm}l@{}}
\toprule
\textbf{Tag} & \textbf{Kind} & \textbf{Meaning} & \textbf{Defined in}\\*
\midrule
\Q{1}--\Q{3} & Question & What the app must answer & Section~\ref{sec:why}\\
\D{1}--\D{9} & Data & A data set and where it comes from & Section~\ref{sec:data}\\
\M{1}--\M{6} & Method & One step that turns data into an output & Sections~\ref{sec:a}--\ref{sec:c}\\
\V{1}--\V{8} & Validation & A check with a pass mark & Section~\ref{sec:valid}\\
\R{1}--\R{8} & Risk & What could go wrong, and the response & Section~\ref{sec:risks}\\
\U{1}--\U{6} & User need & What a user must be able to do & Section~\ref{sec:product}\\
\Sx{1}--\Sx{6} & Software & A component of the app & Section~\ref{sec:arch}\\
\Ox{1}--\Ox{5} & Output & Something the project delivers & Section~\ref{sec:why}\\
\bottomrule
\end{tabular}
\end{center}

A second kind of label says \emph{what backs a number}. It appears next to every value in the tables and on
every number in the app.

\begin{center}
\small\renewcommand{\arraystretch}{1.05}
\begin{tabular}{@{}lp{110mm}@{}}
\toprule
\textbf{Label} & \textbf{Meaning}\\*
\midrule
\Off & A government regulation, official statistic or official open data set.\\
\Mea & A satellite or field measurement.\\
\Tst & Our own method, which passed a test whose pass mark was fixed before the run.\\
\Calc & Arithmetic on backed inputs. It carries their ranges.\\
\Wk & Only one secondary source (news, blog). Shown, but flagged until confirmed.\\
\Scn & An assumption. The app shows it as a slider the user moves, never as a finding.\\
\Mis & Not obtained yet. The app hides the layer until it is.\\
\bottomrule
\end{tabular}
\end{center}

Three kinds of boxes recur. \textcolor{accent}{\textsf{\textbf{In short}}} gives the point of a
section. \textcolor{green}{\textsf{\textbf{In plain words}}} explains an idea without jargon.
\textcolor{warm}{\textsf{\textbf{Honest limit}}} says what a result does not show.

{\sffamily\footnotesize\hypersetup{linkcolor=ink}\setcounter{tocdepth}{1}\setlength{\parskip}{0pt}\makeatletter\renewcommand*\l@section{\@dottedtocline{1}{0em}{2.2em}}\makeatother\tableofcontents}
\clearpage
